% \section{Related Works}

% Our work studies agentic planning under progressively disclosed dual constraints, focusing on how agents handle both world and user constraints during interaction. 
% Table~\ref{tab:comparison-table} provides a trait-level comparison with existing benchmarks, with detailed discussions in Appendix~\ref{app:related-works}.

% A parallel body of work develops methods for agents to better handle constraints during planning.
% For world-side constraints, work tackles several distinct sub-problems.
% On \textit{world state tracking and grounding}, ReflAct \citep{ReflAct} replaces next-action prediction with continuous goal-state reflection, grounding each decision in the current world state to prevent belief–reality misalignment.
% On \textit{constraint violation detection and correction}, \citet{L-ICL} propose Localized In-Context Learning (L-ICL), which identifies the first failing step in a plan trace and injects a targeted minimal correction example, substantially outperforming full-trajectory ICL with far less context.
% On \textit{experience-based world knowledge correction}, XENON \citep{XENON} maintains an adaptive dependency graph and failure-aware action memory that algorithmically revise flawed planning priors from execution experience, enabling robust long-horizon planning in Minecraft even with lightweight models.
% On \textit{long-horizon plan structure and error isolation}, TDP \citep{TDP} decomposes tasks into a DAG of sub-goals so that constraint errors are confined to individual nodes without global propagation, while ALAS \citep{ALAS} coordinates role-specialized stateful agents that apply history-aware local compensation when runtime disruptions occur.
% On \textit{plan generation quality via training}, Plan-and-Act \citep{Plan-and-Act} separates planning from execution into two SFT-trained modules with scalable synthetic data generation, substantially improving structured plan quality on long-horizon tasks.
% On \textit{formal constraint guarantees}, an end-to-end LLM+PDDL framework \citep{end-2-end-pddl} dynamically orchestrates multiple agents to construct and refine PDDL specifications that are handed to a symbolic planner for provably correct constraint satisfaction.
% For user-side constraints, work clusters around two sub-problems.
% On \textit{user preference elicitation}, a line of work trains agents to surface latent preferences through multi-turn interaction: GATE \citep{GATE} introduces a prompting-based generative active elicitation framework using open-ended free-form questions; STaR-GATE \citep{STaR-GATE} iteratively improves the questioner via SFT self-improvement on high-quality elicitation trajectories; and TO-GATE \citep{TO-GATE} further refines this through trajectory optimization, jointly training a clarification resolver and response summarizer.
% On \textit{proactive clarification and personalization in task execution}, Ask-before-Plan \citep{Ask-before-Plan} integrates a SFT-trained clarification agent into the planning pipeline to resolve missing user intent before plan commitment, while PPP \citep{sun2025trainingproactivepersonalizedllm} employs multi-objective RL to jointly optimize productivity, proactivity, and personalization, demonstrating substantial gains over frontier models on software engineering and research tasks.
% A smaller but growing body of work tackles world and user constraints simultaneously, primarily in the travel planning domain. At the prompting and workflow level, MIRROR \citep{MIRROR} is a general tool-use framework whose dual intra-reflection (assessing intended actions before execution) and inter-reflection (adjusting trajectories post-observation) incidentally yields improvements on both commonsense and hard constraint pass rates on TravelPlanner. ATLAS \citep{ATLAS} explicitly addresses this joint challenge through a multi-agent architecture with a dedicated Constraint Manager that identifies both implicit commonsense constraints and explicit user requirements upfront, a Planner–Checker loop for iterative validation, and an adaptive interleaved search that triggers re-planning when any constraint type is violated; it further adapts to evolving user feedback in multi-turn settings. TriFlow \citep{TriFlow} tackles the same problem via a three-stage retrieve–plan–govern pipeline, where a rule–LLM collaboration enforces objective feasibility in the planning stage and bounded iterative refinement in the governance stage aligns outputs with user preferences. At the code-based prompting level, SCOPE \citep{SCOPE} disentangles query-specific reasoning from generic code execution by generating reusable solver functions that encode all constraint types — both commonsense and user hard constraints — as deterministic assertions, achieving state-of-the-art results on TravelPlanner. Finally, HiMAP-Travel \citep{HiMAP-Travel} takes a training-based approach, using GRPO-trained hierarchical agents with a transactional monitor enforcing budget and uniqueness constraints across parallel day-level executors, and a cooperative bargaining protocol that triggers re-planning when any sub-plan violates either objective or user-specified constraints.
% Despite these advances, all of these methods share a common limitation: they are designed for settings where all constraints — both world and user — are disclosed upfront, with no mechanism for discovering constraints progressively through violation-driven interaction. re-plan-Bench directly targets this gap.

